\documentclass[10pt]{article}

\usepackage[utf8]{inputenc}

\usepackage[T1]{fontenc}

\usepackage{amsmath}

\usepackage{amsfonts}

\usepackage{amssymb}

\usepackage[version=4]{mhchem}

\usepackage{stmaryrd}



\begin{document}

в классе любых метрич. пространств. В классе метрич. шространств со счетной базой выполняются $\mathrm{p}$ aв е н с т в а $\mathbf{y}$ ры с о на



\[

\operatorname{dim} X=\operatorname{ind} X=\operatorname{Ind} X \text {. }

\]



В классе любых метрич. пространств выполняется р aв е н с т в о



\[

\operatorname{dim} X=\operatorname{Ind} X

\]



и может быть ind $X=0<\operatorname{Ind} X=1$.



В случае метрич. пространств понятие $n$-мерного пространства следующими двумя способами может быть сведено к понятию нульмерного пространства. Для метрич. пространства $X$ тогда и только тогда $\operatorname{dim} X \leqslant n$, $n=0,1, . .$, когда



a) пространство $X$ может быть представлено в виде не более чем $n+1$ нульмерных слагаемых;



б) существует непрерывное замкнутое отображение кратности $\leqslant n+1$ нульмерного метрич. пространства на пространство $X$.



Для любого подмножества $A$ метрич. пространства $X$ найдется такое подмножество $B \supset A$ типа $G_{\delta}$ в $X$, что $\operatorname{dim} B=\operatorname{dim} A$.



В классе метрич. пространств веса $\leqslant \tau$ п размерности $\leqslant n$ суцествует универсальное (в смысле вложений) пространство. Важную роль в построении теории $\mathrm{P}$. любых метрических (и более общих) пространств сыграла т е о р е м а Д а у к е р а: тогда и только тогда $\operatorname{dim} X \leqslant n$, когда в любое локально конечное открытое покрытие пространства $X$ можно вписать открытое покрытие кратности $\leqslant n+1$.



Одним из наиболее важных вопросов теории Р. является вопрос о соотношениях межкду лебеговой и индуктивными Р. Хотя для произвольного пространства $X$ значения размерностей $\operatorname{dim} X$, ind $X$, Ind $X$, вообще говоря, попарно различны, однако для нек-рых классов пространств, в том или ином смысле близких к метрическим, выполнено, напр., следующее:



a) если пространство $X$ обладает непрерывным замкнутым отображением $f$ размерности $\operatorname{dim} f=0$ на метрич. пространство, то выполняется равенство (3), отсюда следуют равенства (2) для локально бикомпактных групп и их факторпространств;



б) если существует непрерывное замкнутое отображение метрич. пространства на пространство $X$, то выполняются равенства (2).



Еще одно общее условие для выполнения равенства (3) для паракомпакта $X$ выглядит так: $\operatorname{dim} X=n$ и пространство $X$ является образом нульмерного пространства Іри замкнутом отображении кратности $\leqslant n+$ $+1, n=0,1, \ldots$



В случае произвольного пространства $X$ всегда выполняются неравенства $\operatorname{dim} X \leqslant \operatorname{Ind} X$ и ind $X \leqslant \operatorname{Ind} X$, a равенства $\operatorname{dim} X=0$ и Ind $X=0$ равносильны. Для сильно паракомпактного (в частности, бикомпактного или финально компактного) пространства $X$ выполняется неравенство $\operatorname{dim} X \leqslant \operatorname{ind} X$. Для бикомпактов равенства ind $X=1$ и Ind $X=1$ равносильны. Существуют бикомпакты, удовлетворяющие первой аксиоме счетности (и даже совершенно нормальные в предположении континуум-гипотезы), для которых $\operatorname{dim} X=1$, ind $X=n, n=2,3, \ldots$. Построен пример топологич. однородного бикомпакта с $\operatorname{dim} X$ <ind $X$. Для совершенно нормальных бикомпактов всегда ind $X=\operatorname{Ind} X$. Cyществуют бикомпакты даже с первой аксиомой счетности, для к-рых ind $X<\operatorname{Ind} X$. Существует ли такое $m$, что для каждого $n>m$ найдется бикомпакт (метрич. пространство) $X \mathrm{c}$ ind $X=m$, Ind $X=n,-$ неизвестно (1983).



В случае неметризуемых пространств $\mathrm{P}$. может не только не быть монотонной, но и обладает другими патологич. свойствами. Для любого $n=2,3, \ldots$. построен пример такого бикомпакта $\theta_{n}$, что любое замкнутое подмножество его имеет P. или 0 или $n=\operatorname{dim} \theta_{n}$. Аналогичный пример в случає индуктивных $\mathrm{P}$. невозможен. Построен также для любого $n=1,2, \ldots$.пример такого бикомпакта $\Phi_{n}$, что любое разбивающее этот бикомпакт замкнутое множество имеет размерность $n=$ $=\operatorname{dim} \Phi_{n}$. Таким образом, подход к определению Р. в случае неметризуемого пространства в принципе отличен от индуктивного подхода А. Пуанкаре, основанного на разбиении пространства подпространствами меньшего числа измерений. Бикомпакты $\Phi_{n}$ имеют непосредственное отношение к следующему утверждению: в любом n-мерном бикомпакте содержится $n$-мерное канторово многообразие.



Подмножество $n$-мерного евклидова пространства $E^{n}$ тогда и только тогда $n$-мерно, когда оно содержит внутренние относительно $E^{n}$ точки. Компакт имеет размерность $\leqslant n$ тогда и только тогда, когда он обладает отображением Р. нуль в $E^{n}$, и, таким образом, с точностью до нульмерных отображений $n$-мерные компакты не отлиqимы от ограниченных замкнутых, содержащих внутренние (относительно $E^{n}$ ) точки подмножеств $E^{n}$.



См. также Размерности теория.



Лuт.: [1] А ле к с а нд ро в П. С., Па сын ко в Б. А., Введение в теорию размерности, М., 1973; [2] Г у р е в и ч В., В о л м ә н Г., Теория размерности, пер. с англ., М., 1948; [3] У р ы с о н П. С., Труды по топологии и другим областям матеРАЗМЕЩЕНИЕ с П о в т о р е н и я М И И элементов по $n$ - конечная последовательность $\bar{a}=$ $=\left(a_{i_{1}}, a_{i_{2}}, \ldots, a_{i_{n}}\right)$ элементов нек-рого множества $A=$ $=\left\{a_{1}, a_{2}, \ldots, a_{m}\right\}$. Если все члены $\bar{a}$ различны, то $\bar{a}$ наз. Р. бе з по в то р е н йй. Число всех возможных Р. с повторениями из $m$ по $n$ равно $m^{n}$, а без повторений $(m)_{n}==m(m-1) \ldots(m-n-1)$.



Р. можно рассматривать как функцию $\varphi$, заданную на $Z_{n}=\{1,2, \ldots, n\}$ и принимающую значения из $A$ : $\varphi(k=) a_{i_{k}}, k=1,2, \ldots, n$. Элементы $\boldsymbol{A}$ принлто называть я ч е й к а м и (или у р н а м и), а әлементы $Z_{n}-$ ч ас т и ц а м и (или ші а р а м и ); $\varphi$ определяөт заполнение различных ячеек различными частицами. Если речь идет о неразличимых частицах или ячейках, то подразумевается, что рассматриваются к л а с с Ы Р. Так, если все частицы одинаковы, то два Р., определяемые соответственно функциями $\varphi$ и $\psi$, относятся к одному классу, если найдегся подстановка $\sigma$ множества $Z_{n}$ такая, что $\varphi(\sigma(k))=\psi(k)$ для всех $k \in Z_{n}$. В этом случае число таких классов, или, как говорят, число размещений $n$ одинаковых частиц по $m$ различным ячейкам, есть число сочетаний с повторениями из $n$ по $m$.



Если говорят, что все ячейки одинаковы, то имеют в виду, что Р. разбиваются на классы так, что два Р., определяемые функциями $\varphi$ и $\psi$ соответственно, относятся к одному классу, если существует подстановка $\tau$ множества $A$, при к-рой $\tau(\varphi(k))=\psi(k)$ для всех $k \in Z_{n}$. В этом случае число размещений $n$ различных частиц по $m$ одинаковым ячейкам, т. е. число классов, равно $\sum_{k=1}^{m} S(n, k)$, где $S(n, k)$ - ч и с л а С т и рл и нг а II р о д а:



$S(n, k)=S(n-1, k-1)+k S(n-1, k), S(0,0)=1$, $S(n, k)=0$ при $k>n \quad$ и $k=0, n>0$.



Если не различать как частицы, так и ячөйки, то получают размецение $n$ одинаковых частиц по $m$ одинаковым ячейкам; число таких Р. равно $\sum_{k=1}^{n} p_{n}(k)$, где $p_{n}(k)$ - число разбиений $n$ на $k$ натуральных слагаемых.





\end{document}